\documentclass[../proyecto.tex]{subfiles}
\graphicspath{{\subfix{../images/}}}
\begin{document}

\chapter{Resumen}

Este proyecto de tesis es una investigación artística que propone el concepto de popusintesíntesis, esto es, la síntesis popular de sintetizadores, en particular los sonoros.

La misma palabra popusintesíntesis es una palabra inventada, o simplemente palabra, que hereda la práctica del artista y empresario Peter Blasser, quien define su práctica artística como sintesíntesis, dado que él se dedica a la síntesis de sintetizadores.

En esta tesis se proponen conceptualizaciones de sintetizadores sonoros, que son capaces de producir sonido, a partir de la interconexión de distintos módulos.

Un aporte de esta tesis es la creación de materiales remezclables, que permitan que emerjan nuevos sintetizadores a partir de este trabajo de creación.

También una novedad que tiene este proyecto de tesis es que se propone una metodología de creación que se define como popular, esto es, se hace de forma pública, sostenida en el tiempo, con la intención de que grandes multitudes puedan tener acceso a la educación, fabricación y consumo de sintetizadores sonoros.

% máximo 1 página

\end{document}